\section{Overview and Principles}
All statistical analysis is performed at the \emph{clip level}, where each clip is evaluated under
multiple paired conditions (e.g., preprocessing A vs B) on the \emph{same} audio content. This
paired design reduces variance and enables stronger inference about the effect of overlap, SNR,
preprocessing, and model choice. :contentReference[oaicite:1]{index=1}

Unless otherwise stated, all results are reported as mean $\pm$ 95\% confidence interval (CI)
computed via bootstrap resampling over clips within a condition (Section~\ref{sec:bootstrap_ci}).

\section{Response Variables}
For each experimental unit (clip $\times$ preprocessing $\times$ embedding setting $\times$ ASR model),
the following response variables are computed:
\begin{itemize}
  \item $y_{WER}$: Word Error Rate (WER) over the full clip reference.
  \item $y_{oWER}$: overlap-only WER, computed only on words that occur within overlapped regions.
  \item $y_{cpWER}$: cpWER for multi-speaker reference scoring when applicable.
\end{itemize}
Overlap-only scoring uses the overlap mask stored in synthetic mixtures, and overlap annotations/estimates
for real audio (Chapter~3--4). :contentReference[oaicite:2]{index=2}

\section{Bootstrap Confidence Intervals}
\label{sec:bootstrap_ci}
For each condition cell (defined by overlap ratio, $K$, SNR, noise type, preprocessing, and ASR model),
a stratified bootstrap over clips is used to estimate uncertainty. Let $\{y_i\}_{i=1}^n$ be clip-level
scores within the cell. We sample with replacement $n$ clips for $B$ resamples (e.g., $B=10{,}000$),
compute the mean score for each resample, and take the 2.5th and 97.5th percentiles to form a 95\% CI.

\section{Paired Comparisons (Preprocessing Effects)}
To answer RQ4 and RQ5, preprocessing pipelines are compared \emph{paired} on identical clips:
\[
d_i = y_i^{(A)} - y_i^{(B)} ,
\]
where $y_i^{(A)}$ and $y_i^{(B)}$ are clip-level scores under pipelines A and B, respectively.
We report:
\begin{itemize}
  \item mean difference $\bar{d}$ and 95\% CI via paired bootstrap;
  \item a non-parametric paired test (Wilcoxon signed-rank) as a robustness check.
\end{itemize}
Because multiple preprocessing comparisons are performed across many conditions and models,
$p$-values are adjusted using Holm--Bonferroni correction.

\section{Main Effects and Interactions (Model-Based Analysis)}
To quantify interactions such as preprocessing $\times$ overlap ratio and model $\times$ overlap ratio,
a regression-style analysis is applied on clip-level scores. Let the design factors be:
overlap ratio (OR), max concurrent speakers ($K$), SNR, noise type, preprocessing pipeline, ASR model,
and speaker embedding availability (when relevant). The analysis focuses on:
\begin{itemize}
  \item \textbf{main effects}: OR, $K$, SNR, preprocessing, model;
  \item \textbf{key interactions}: OR$\times$preprocessing, OR$\times$model, preprocessing$\times$model.
\end{itemize}

A mixed-effects model is used when feasible to account for repeated measurements per clip:
\[
y_{i} = \beta_0 + \beta^\top X_i + u_{\text{clip}(i)} + \epsilon_i,
\]
where $X_i$ contains fixed effects and interaction terms, $u_{\text{clip}(i)}$ is a random intercept for clip,
and $\epsilon_i$ is residual error.

If mixed-effects estimation is unstable for WER distributions, the primary inference falls back to:
(i) paired bootstrap comparisons for effect sizes, and
(ii) condition-wise aggregated analysis with robust standard errors.

\section{Reporting of Effect Sizes}
Beyond statistical significance, the following are reported:
\begin{itemize}
  \item absolute WER change: $\Delta WER$;
  \item relative WER reduction: $(WER_B - WER_A)/WER_B$;
  \item overlap-only improvements $\Delta oWER$ as the primary overlap-focused outcome.
\end{itemize}

\section{Transferability: Synthetic vs Real (RQ3)}
To assess transferability, model and preprocessing rankings (by WER and overlap-only WER)
are computed on synthetic mixtures and real datasets separately. Agreement is quantified using
rank correlation (Spearman) across systems and key operating points (e.g., OR $\ge 20\%$, SNR $\le 0$ dB).